% GENERATED FILE. Edit canonical lessons, then run npm run generate:books.
% canonical-source-hash: fnv1a64:02624209eb4bad67
% canonical-lessons: JA-C247-homeru, JA-C247-shikaru, JA-C247-hazukashii, JA-C247-komaru, JA-C247-yuuki, JA-R247-review-242-247

\chapter{Praise and Scolding}
\label{ch:ja-praise-and-scolding}
\hlchaptermodality{247}

\begin{chapteropening}
\textbf{By the end of this chapter:} \emph{I can ask who praises and who scolds.}

Complete the last lesson of chapter 247: I can ask who praises and who scolds.
\end{chapteropening}

\section[\texorpdfstring{\ja{ほめる} (homeru)}{ほめる (homeru)}]{\ja{ほめる} (homeru) — to praise}

\label{lesson:JA-C247-homeru}



\begin{quote}
Before the new one: say the Japanese for to forgive, then the Japanese for a pity.
\end{quote}

\subsection*{You'll want to know: \ja{ほめる}}
\textbf{\ja{ほめる}} — \emph{homeru} — \textquotedblleft{}to praise\textquotedblright{}.

\begin{quote}
\textbf{\ja{せんせいは} \ja{よく} \ja{がくせいを} \ja{ほめますか。}} — \emph{sensei wa yoku gakusei o homemasu ka} — \textquotedblleft{}Does the teacher often praise the students?\textquotedblright{}
\end{quote}

Polite \textbf{\ja{ほめます}}, \emph{homemasu}; an ichidan verb.

\begin{cousinweb}
\textbf{\ja{ほめる}} (\emph{homeru}) is a native Japanese word (wago), inherited rather than borrowed from Chinese; its kanji was attached later as a kun'yomi (JMdict, KANJIDIC2).
\end{cousinweb}

\begin{grammarlens}[title={What you've built}]
One more word for everyday life at A2.
\end{grammarlens}

\subsection*{Your turn}
\begin{itemize}
\raggedright
  \item \emph{Say it:} \emph{homeru}
  \item \emph{Say it:} \emph{homeru}, once more
  \item \emph{Say it:} \emph{sensei wa yoku gakusei o homemasu ka} — Does the teacher often praise the students?
  \item \emph{Recall:} say the Japanese for to forgive, then the Japanese for a pity, then say \emph{homeru} again
\end{itemize}

\begin{tcolorbox}[breakable,colback=teal!4,colframe=teal!35!black,title={Before you move on}]
What does \ja{ほめる} mean? (\textquotedblleft{}To praise\textquotedblright{}.) Say it once more.
\end{tcolorbox}
\section[\texorpdfstring{\ja{しかる} (shikaru)}{しかる (shikaru)}]{\ja{しかる} (shikaru) — to scold}

\label{lesson:JA-C247-shikaru}



\begin{quote}
Before the new one: say the Japanese for a pity, then the Japanese for to praise.
\end{quote}

\subsection*{You'll want to know: \ja{しかる}}
\textbf{\ja{しかる}} — \emph{shikaru} — \textquotedblleft{}to scold\textquotedblright{}.

\begin{quote}
\textbf{\ja{そふは} \ja{こどもを} \ja{しかりません。}} — \emph{sofu wa kodomo o shikarimasen} — \textquotedblleft{}My grandfather does not scold the children.\textquotedblright{}
\end{quote}

Polite \textbf{\ja{しかります}}, \emph{shikarimasu}; a godan verb.

\begin{cousinweb}
\textbf{\ja{しかる}} (\emph{shikaru}) is a native Japanese word (wago), inherited rather than borrowed from Chinese; its kanji was attached later as a kun'yomi, and the word is usually written in kana (JMdict, KANJIDIC2).
\end{cousinweb}

\begin{grammarlens}[title={What you've built}]
One more word for everyday life at A2.
\end{grammarlens}

\subsection*{Your turn}
\begin{itemize}
\raggedright
  \item \emph{Say it:} \emph{shikaru}
  \item \emph{Say it:} \emph{shikaru}, once more
  \item \emph{Say it:} \emph{sofu wa kodomo o shikarimasen} — My grandfather does not scold the children.
  \item \emph{Recall:} say the Japanese for a pity, then the Japanese for to praise, then say \emph{shikaru} again
\end{itemize}

\begin{tcolorbox}[breakable,colback=teal!4,colframe=teal!35!black,title={Before you move on}]
What does \ja{しかる} mean? (\textquotedblleft{}To scold\textquotedblright{}.) Say it once more.
\end{tcolorbox}
\section[\texorpdfstring{\ja{はずかしい} (hazukashii)}{はずかしい (hazukashii)}]{\ja{はずかしい} (hazukashii) — embarrassed; shy}

\label{lesson:JA-C247-hazukashii}



\begin{quote}
Before the new one: say the Japanese for to praise, then the Japanese for to scold.
\end{quote}

\subsection*{You'll want to know: \ja{はずかしい}}
\textbf{\ja{はずかしい}} — \emph{hazukashii} — \textquotedblleft{}embarrassed; shy\textquotedblright{}.

\begin{quote}
\textbf{\ja{いま} \ja{はずかしいですか。}} — \emph{ima hazukashii desu ka} — \textquotedblleft{}Are you embarrassed now?\textquotedblright{}
\end{quote}

\begin{cousinweb}
\textbf{\ja{はずかしい}} (\emph{hazukashii}) is a native Japanese word (wago), inherited rather than borrowed from Chinese; its kanji was attached later as a kun'yomi (JMdict, KANJIDIC2).
\end{cousinweb}

\begin{grammarlens}[title={What you've built}]
One more word for everyday life at A2.
\end{grammarlens}

\subsection*{Your turn}
\begin{itemize}
\raggedright
  \item \emph{Say it:} \emph{hazukashii}
  \item \emph{Say it:} \emph{hazukashii}, once more
  \item \emph{Say it:} \emph{ima hazukashii desu ka} — Are you embarrassed now?
  \item \emph{Recall:} say the Japanese for to praise, then the Japanese for to scold, then say \emph{hazukashii} again
\end{itemize}

\begin{tcolorbox}[breakable,colback=teal!4,colframe=teal!35!black,title={Before you move on}]
What does \ja{はずかしい} mean? (\textquotedblleft{}Embarrassed; shy\textquotedblright{}.) Say it once more.
\end{tcolorbox}
\section[\texorpdfstring{\ja{こまる} (komaru)}{こまる (komaru)}]{\ja{こまる} (komaru) — to be in trouble; to be at a loss}

\label{lesson:JA-C247-komaru}



\begin{quote}
Before the new one: say the Japanese for to scold, then the Japanese for embarrassed.
\end{quote}

\subsection*{You'll want to know: \ja{こまる}}
\textbf{\ja{こまる}} — \emph{komaru} — \textquotedblleft{}to be in trouble; to be at a loss\textquotedblright{}.

\begin{quote}
\textbf{\ja{あには} \ja{しごとで} \ja{こまりますか。}} — \emph{ani wa shigoto de komarimasu ka} — \textquotedblleft{}Does my elder brother have trouble at work?\textquotedblright{}
\end{quote}

Polite \textbf{\ja{こまります}}, \emph{komarimasu}; a godan verb.

\begin{cousinweb}
\textbf{\ja{こまる}} (\emph{komaru}) is a native Japanese word (wago), inherited rather than borrowed from Chinese; its kanji was attached later as a kun'yomi (JMdict, KANJIDIC2).
\end{cousinweb}

\begin{grammarlens}[title={What you've built}]
One more word for everyday life at A2.
\end{grammarlens}

\subsection*{Your turn}
\begin{itemize}
\raggedright
  \item \emph{Say it:} \emph{komaru}
  \item \emph{Say it:} \emph{komaru}, once more
  \item \emph{Say it:} \emph{ani wa shigoto de komarimasu ka} — Does my elder brother have trouble at work?
  \item \emph{Recall:} say the Japanese for to scold, then the Japanese for embarrassed, then say \emph{komaru} again
\end{itemize}

\begin{tcolorbox}[breakable,colback=teal!4,colframe=teal!35!black,title={Before you move on}]
What does \ja{こまる} mean? (\textquotedblleft{}To be in trouble; to be at a loss\textquotedblright{}.) Say it once more.
\end{tcolorbox}
\section[\texorpdfstring{\ja{ゆうき} (yūki)}{ゆうき (yūki)}]{\ja{ゆうき} (yūki) — courage}

\label{lesson:JA-C247-yuuki}



\begin{quote}
Before the new one: say the Japanese for embarrassed, then the Japanese for to be in trouble.
\end{quote}

\subsection*{You'll want to know: \ja{ゆうき}}
\textbf{\ja{ゆうき}} — \emph{yūki} — \textquotedblleft{}courage\textquotedblright{}.

\begin{quote}
\textbf{\ja{かれは} \ja{ゆうきが} \ja{ありますか。}} — \emph{kare wa yūki ga arimasu ka} — \textquotedblleft{}Does he have courage?\textquotedblright{}
\end{quote}

\begin{cousinweb}
\textbf{\ja{ゆうき}} (\emph{yūki}) is Sino-Japanese (kango): two Chinese characters read with their on'yomi, readings borrowed from Chinese. They are \emph{yū}, \textquotedblleft{}brave\textquotedblright{} (Middle Chinese \emph{yowngX}) and \emph{ki}, \textquotedblleft{}spirit, air\textquotedblright{} (\emph{khj+jH}) (KANJIDIC2; Middle Chinese in Baxter's transcription).
\end{cousinweb}

\begin{grammarlens}[title={What you've built}]
One more word for everyday life at A2.
\end{grammarlens}

\subsection*{Your turn}
\begin{itemize}
\raggedright
  \item \emph{Say it:} \emph{yūki}
  \item \emph{Say it:} \emph{yūki}, once more
  \item \emph{Say it:} \emph{kare wa yūki ga arimasu ka} — Does he have courage?
  \item \emph{Recall:} say the Japanese for embarrassed, then the Japanese for to be in trouble, then say \emph{yūki} again
\end{itemize}

\begin{tcolorbox}[breakable,colback=teal!4,colframe=teal!35!black,title={Before you move on}]
What does \ja{ゆうき} mean? (\textquotedblleft{}Courage\textquotedblright{}.) Say it once more.
\end{tcolorbox}
\section[(dialogue)]{Words from chapters 242–247 — a review}

\label{lesson:JA-R247-review-242-247}



\begin{quote}
Say the Japanese for he and she.
\end{quote}

\subsection*{Your turn}
Say these aloud:

\begin{itemize}
\raggedright
  \item he, she, a sweetheart, a close friend, to go out with
  \item parents, my grandfather, my grandmother, siblings, to bring up
  \item an acquaintance, to come across, to part, to welcome, marriage
  \item glad, lonely, to be delighted, to laugh, a smile
  \item to get angry, a quarrel, to apologize, to forgive, a pity
  \item to praise, to scold, embarrassed, to be in trouble, courage
\end{itemize}

\begin{tcolorbox}[breakable,colback=teal!4,colframe=teal!35!black,title={Before you move on}]
He? (\textbf{\ja{かれ}}, \emph{kare}.) She? (\textbf{\ja{かのじょ}}, \emph{kanojo}.) A sweetheart? (\textbf{\ja{こいびと}}, \emph{koibito}.) A close friend? (\textbf{\ja{しんゆう}}, \emph{shin'yū}.) To go out with? (\textbf{\ja{つきあう}}, \emph{tsukiau}.) Parents? (\textbf{\ja{りょうしん}}, \emph{ryōshin}.) My grandfather? (\textbf{\ja{そふ}}, \emph{sofu}.) My grandmother? (\textbf{\ja{そぼ}}, \emph{sobo}.) Siblings? (\textbf{\ja{きょうだい}}, \emph{kyōdai}.) To bring up? (\textbf{\ja{そだてる}}, \emph{sodateru}.) An acquaintance? (\textbf{\ja{しりあい}}, \emph{shiriai}.) To come across? (\textbf{\ja{であう}}, \emph{deau}.) To part? (\textbf{\ja{わかれる}}, \emph{wakareru}.) To welcome? (\textbf{\ja{むかえる}}, \emph{mukaeru}.) Marriage? (\textbf{\ja{けっこん}}, \emph{kekkon}.) Glad? (\textbf{\ja{うれしい}}, \emph{ureshii}.) Lonely? (\textbf{\ja{さびしい}}, \emph{sabishii}.) To be delighted? (\textbf{\ja{よろこぶ}}, \emph{yorokobu}.) To laugh? (\textbf{\ja{わらう}}, \emph{warau}.) A smile? (\textbf{\ja{えがお}}, \emph{egao}.) To get angry? (\textbf{\ja{おこる}}, \emph{okoru}.) A quarrel? (\textbf{\ja{けんか}}, \emph{kenka}.) To apologize? (\textbf{\ja{あやまる}}, \emph{ayamaru}.) To forgive? (\textbf{\ja{ゆるす}}, \emph{yurusu}.) A pity? (\textbf{\ja{ざんねん}}, \emph{zannen}.) To praise? (\textbf{\ja{ほめる}}, \emph{homeru}.) To scold? (\textbf{\ja{しかる}}, \emph{shikaru}.) Embarrassed? (\textbf{\ja{はずかしい}}, \emph{hazukashii}.) To be in trouble? (\textbf{\ja{こまる}}, \emph{komaru}.) Courage? (\textbf{\ja{ゆうき}}, \emph{yūki}.)
\end{tcolorbox}
